\section{Introduction}
\label{sec:intro}

Overhead aerial imagery feeds segmentation pipelines that are deployed across geographic regions. Shadow masks are a core input to several of these pipelines: solar irradiance estimation uses them to separate shaded from sunlit surfaces \cite{batchu2024satellite}, urban heat mapping uses them to correct surface temperature readings \cite{dashadebench, ukkusuri2024heatwaves}, and building footprint extraction uses them to resolve occlusion boundaries \cite{seo}. These pipelines are routinely trained on data from one region and deployed in another~\cite{mittal2025overcoming, verma2025spatial}.

Geographic shift is a known failure mode in segmentation models across domains. Medical imaging models underperform across hospitals due to scanner and protocol differences~\cite{ogier2022flamby}, and aerial land-cover models show similar drops on unseen regions~\cite{wang2021loveda, gong2025crossearth}. For overhead shadow detection, no comparable evaluation infrastructure exists. AISD, the largest human-verified overhead shadow dataset, contains 514 images and defines no train-test split across cities~\cite{aisd}. S-EO scales further but relies on automatically generated masks with documented alignment artifacts, and again reports no geographic holdout~\cite{seo}. Other widely used shadow benchmarks (SBU, ISTD) target ground-level imagery and do not match overhead geometry~\cite{sbu, istd}. Table~\ref{tab:dataset_comparison} summarizes the landscape. Without a benchmark that pairs human-verified labels with a defined geographic protocol, transfer failure in overhead shadow detection cannot be measured, diagnosed, or addressed.

\begin{table}[t]
  \centering
  \caption{
    \textbf{Comparison of shadow detection datasets.}
    No existing overhead dataset combines human-verified labels, multiple cities, multiple native resolutions, and a defined geographic transfer protocol.
    \shadow{} fills this gap.
    ISTD is included for context only; ground-level geometry is not a substitute for overhead evaluation.
    ``Transfer protocol'' = a defined geographic train-test split with formal evaluation guidelines.
    $\checkmark$~=~present; \texttimes{}~=~absent; n/r~=~not reported.
  }
  \label{tab:dataset_comparison}
  \small
  \setlength{\tabcolsep}{4pt}
  \begin{tabular}{@{}lllcccc@{}}
    \toprule
    Dataset & Domain & Scale & Human QC & Cities & Multi-res & Transfer protocol \\
    \midrule
    ISTD~\cite{istd}                & Ground-level             & 1{,}870     & $\checkmark$         & --          & \texttimes & \texttimes \\
    AISD~\cite{aisd}                & Aerial 0.3\,m   & 514         & $\checkmark$         & 5           & \texttimes & \texttimes \\
    Airborne-Shadow~\cite{azimi2023airborne} & Aerial & n/r         & $\checkmark$         & DE only     & \texttimes & \texttimes \\
    S-EO~\cite{seo}                 & Satellite 0.3\,m & $\sim$19{,}000 & \texttimes\ (auto) & 3 (US)      & \texttimes & \texttimes \\
    \midrule
    \shadow{} (ours)                & Aerial 0.3/0.6\,m & 4{,}500   & $\checkmark$         & 3 (US)      & $\checkmark$ & $\checkmark$\ \loco \\
    \bottomrule
  \end{tabular}
\end{table}
\begin{figure}[t]
  \centering
  \includegraphics[width=\linewidth]{fig1_teaser.png}
  \caption{
    \textbf{\shadow and the geographic transfer gap.}
    \emph{Left:} paired test patches from three U.S.\ cities at two native NAIP
    resolutions ($0.3$\,m/px, $0.6$\,m/px); each column shares ground
    coordinates between rows, and the white box on each $0.6$\,m/px tile
    marks the paired $0.3$\,m/px footprint. Shadow masks are human-verified
    across six categories (bottom legend).
    \emph{Right:} two-class boundary-tolerant mIoU under vanilla training on the high-resolution cells, with each city held out in turn and results averaged over the three rotations (inset). Open marker $=$ within-city upper bound; filled marker $=$ LOCO held-out city; $\Delta =$ mean gap; error bars are $95\%$ CIs. Methods: Table~\ref{tab:loco-main}.
}
  \label{fig:teaser}
\end{figure}


To bridge this gap, we introduce \shadow, a benchmark for measuring geographic transfer in overhead shadow detection. It contains 4{,}500 human-verified shadow masks across three U.S. cities (Chicago, Miami, Phoenix) at two native sensor resolutions (0.3 m and 0.6 m), with a leave-one-city-out (LOCO) protocol that trains on two cities and tests on the third. 

Figure~\ref{fig:teaser} shows the size of the gap this protocol exposes: under vanilla training on the high-resolution cells, holding out the test city costs between $3.5$ mIoU (CrossEarth) and $7.8$ mIoU (OGLANet), and existing domain-adaptation methods help on some architecture-city pairs and fail catastrophically on others. The major \textbf{contributions} of this work are threefold:

\begin{enumerate}
    \item We release \shadow{}\footnote{Data: \url{https://huggingface.co/datasets/shagunmittal/ShadowTransfer} (DOI \href{https://doi.org/10.57967/hf/11483}{10.57967/hf/11483}); code: \url{https://github.com/shagun8/shadowtransfer}; project page: \url{https://shagun8.github.io/shadowtransfer/}; full datasheet in App.~\ref{app:datasheet}.} with its LOCO benchmarking and statistical reporting. 
    \item We provide a three-probe evaluation framework that characterizes \emph{how} models fail under geographic transfer: \textit{where} in the pipeline the shift acts, \textit{how} the prediction distribution distorts, and \textit{which} architectural components carry the failure
    \item We present \sib, a case-study module designed from the framework's findings, with the smallest mean gap to the within-city upper bound of the ten methods we evaluate, level with FDA. Together, the benchmark and framework give practitioners a way to measure geographic transfer risk before deployment and to diagnose where in a model that risk concentrates.
\end{enumerate}

